\documentclass[11pt]{article}

% --------------------------------------------------
% Packages
% --------------------------------------------------
\usepackage{booktabs}   % for \toprule, \midrule, \bottomrule
\usepackage{caption}    % for \captionof
\usepackage{geometry}
\usepackage{setspace}
\usepackage{amsmath}    % for equations
\usepackage{graphicx}

\geometry{margin=1in}
\setstretch{1.1}

% --------------------------------------------------
% Title
% --------------------------------------------------
\title{Replication Summary\\
\large Szymanowska et al. (2014)}
\author{Garbuzov and Mitchell}
\date{\today}

% --------------------------------------------------
% Document
% --------------------------------------------------
\begin{document}

\maketitle

\section{Data Construction and Descriptives}

In this paper, we replicate Tables I and II from Szymanowska et al. (2014), which deconstruct the spot and term premia for various commodity sectors.
The data is analyzed from March 1986 to December 2010, and is sourced from the Wharton Research Data Services (WRDS) database. 
Succeses in this project were a better understanding of commodities premia, deployment of ideas learned from class such as juypter notebooks and cruft, and succesfully building a github repo from start to finish!
Challenges in this project were replication of the paper, creating tables within a Latex document, and finding commodities futures contracts information.

\subsection{Commodity Panel Sample (Head)}

Below is the head of the cleaned commodity panel used in our analysis. 
In the paper, contracts are extracted at 2 month intervals, which equate to an interval of n = 1. The original paper then takes the closest 2 month contract to be the spot price. 
Note, we chose to extract values on the last day of every other month, as the date of extraction was not specified by the paper. Beyond that, the specific commodities codes to use were not specified.

\begin{center}
\resizebox{0.9\textwidth}{!}{%
  \input{../_output/commodity_panel_head.tex}%
}
\end{center}

\subsection{Sector Classification List}

For this paper, seven commodities sectors are discussed. Those are Energy, Grains, Industrial Materials, Meats, Metals, Oilseeds, and Softs.



\section{Return Construction}

Returns are computed from bimonthly futures prices at different
maturities. We follow the paper’s notation and define Short Roll and
Excess Holding returns as functions of log futures prices.

\subsection{Short Roll Returns}

Let \(F_{t,n}\) denote the futures price at date \(t\) with \(n\)
bimonthly periods to maturity, and let \(f_{t,n} = \ln F_{t,n}\) be its
log price. The Short Roll (SR) return over one bimonthly period for
maturity \(n\) is

\[
r^{SR}_{t+1,n}
  = f_{t+1,n} - f_{t,n+1}.
\]

When \(n=1\), this represents rolling from the near contract into the
next contract while keeping maturity approximately constant. Note that the roll over date is implied to be 4 to 6 weeks earlier than traditional rolling over to deal with liquidity issues.

\subsection{Excess Holding Returns}

The Excess Holding (EH) return over \(n\) bimonthly periods is defined
as the \(n\)-period futures return in excess of the corresponding
short-roll strategy. Denote the \(n\)-period log futures return by

\[
r^{F}_{t \rightarrow t+n}
  = f_{t+n,n} - f_{t,n+n}.
\]

Then the Excess Holding return for horizon \(n\) is

\[
r^{EH}_{t \rightarrow t+n}
  = r^{F}_{t \rightarrow t+n}
    - \sum_{j=0}^{n-1} r^{SR}_{t+j,1}.
\]

By construction, EH isolates the payoff to holding the futures position
to maturity relative to rolling short at the front of the curve. This method allows us to break apart a futures contract into the spot and term premia components.

\section{Summary Statistics Tables}

Breaking apart futures contracts by shorting the closest contract and longing the further away contract, we calculcate the term premia associated with that contract. From our dataset (March 1986--December 2010), we then compute annualized mean returns, standard deviations, and t-statistics for Short Roll and Excess Holding returns across commodity sectors, replicating the main results of the paper.
Our findings are quite different from the original paper, which is likely due to differences in data sources and return construction details that were not fully specified in the original paper.

\subsection{Main Table 1}

Table~\ref{tab:table1} reports annualized mean returns, standard
deviations, and Newey--West \(t\)-statistics for Short Roll and Excess
Holding returns across commodity sectors.

\begin{table}[htbp]
\centering
\caption{Summary Statistics}
\label{tab:table1}
\footnotesize
\begin{tabular}{lrrrrrrrrrrrr}
\toprule
 & \multicolumn{4}{c}{Mean} & \multicolumn{4}{c}{Std Dev} & \multicolumn{4}{c}{$t$-stat} \\
 & $n=1$ & $n=2$ & $n=3$ & $n=4$ & $n=1$ & $n=2$ & $n=3$ & $n=4$ & $n=1$ & $n=2$ & $n=3$ & $n=4$ \\
\midrule
\textit{Short Roll} \\
Energy & 8.88% & 8.86% & 8.82% & — & 34.61% & 33.05% & 31.62% & — & (1.32) & (1.37) & (1.41) & — \\
Meats & 0.39% & 3.65% & 2.81% & — & 12.96% & 11.71% & 11.01% & — & (0.13) & (1.33) & (1.19) & — \\
Metals & 4.64% & 4.15% & 3.62% & — & 17.47% & 17.60% & 17.51% & — & (1.23) & (1.09) & (0.94) & — \\
Grains & -7.16% & -2.35% & -3.94% & — & 19.02% & 18.73% & 17.48% & — & (-1.55) & (-0.48) & (-0.85) & — \\
Oilseeds & 1.24% & 2.67% & 4.02% & — & 23.66% & 23.20% & 22.14% & — & (0.29) & (0.61) & (0.93) & — \\
Softs & -7.82% & -7.15% & -7.00% & — & 19.02% & 18.11% & 17.61% & — & (-2.09) & (-1.95) & (-1.91) & — \\
Ind materials & -5.84% & -1.87% & -0.90% & — & 19.80% & 17.28% & 15.40% & — & (-1.28) & (-0.48) & (-0.26) & — \\
EW & -1.62% & 0.53% & 0.23% & — & 11.68% & 11.48% & 11.06% & — & (-0.59) & (0.19) & (0.09) & — \\
\midrule
\textit{Excess Holding} \\
Energy & — & -0.01% & -0.02% & — & — & 2.30% & 3.66% & — & — & (-0.04) & (-0.05) & — \\
Meats & — & 1.63% & 1.95% & — & — & 3.27% & 4.30% & — & — & (3.84) & (3.08) & — \\
Metals & — & -0.37% & -0.65% & — & — & 0.64% & 0.89% & — & — & (-4.22) & (-4.44) & — \\
Grains & — & 2.40% & 2.64% & — & — & 2.91% & 3.87% & — & — & (5.04) & (3.96) & — \\
Oilseeds & — & 0.71% & 1.41% & — & — & 2.00% & 3.42% & — & — & (2.41) & (2.89) & — \\
Softs & — & 0.33% & 0.53% & — & — & 1.69% & 2.56% & — & — & (1.21) & (1.20) & — \\
Ind materials & — & 1.99% & 3.06% & — & — & 4.02% & 5.50% & — & — & (2.70) & (2.87) & — \\
EW & — & 1.06% & 1.35% & — & — & 1.22% & 1.80% & — & — & (4.81) & (4.08) & — \\
\bottomrule
\end{tabular}
\end{table}

\newpage
\subsection{Extended Table 1}

The extended version of Table~1 below uses the extended
sample up until the present-day and applies the same return definitions and sector groupings.

\input{../_output/table1_extended.tex}

\subsection{Additional Summary Statistics}

This table reports additional summary statistics for the commodity
panel, beyond those in Table~1. We compute statistics across the same length as the paper, March 1986--December 2010.
The mean basis indicates to us whether a sector is on average in backwardation or contango. Energy and metals exhibit near-zero or negative mean basis, 
while meats and grains are predominantly in contango. Note the std, minimum and maximum basis values as well, which indicate a high level of volatility in many commodities. Skew and kurtosis further help us understand the general shape of these returns.

\input{../_output/summary_stats.tex}

\section{Sector-Level Short Roll Premia}

This chart shows annualized Short Roll mean returns
by sector and maturity, using the sector rows from Table~1. Similair to the summary statistics above, SR premia indicate whether a sector is on average in backwardation or contango. We've graphed spot premia for maturities \(n=1,2,3,4\) to show how the premia evolve as we move further out along the futures curve.
\begin{figure}[h]
  \centering
  \includegraphics[width=0.8\textwidth]{../_output/fig_sr_sector_bar.png}
  \caption{Sector-level Short Roll (SR) annualized mean returns by
  maturity $n=1,2,3,4$.}
  \label{fig:sr_sector_bar}
\end{figure}


\end{document}
